\def\uadectitulo{Redacción indirecta}
\documentclass[12pt,letterpaper,oneside]{article}
\usepackage[T1]{fontenc}
\usepackage[utf8]{inputenc}
\usepackage[spanish,es-nodecimaldot]{babel}
\usepackage{lmodern}
\usepackage[margin=2.5cm]{geometry}
\usepackage{amsmath,amssymb,graphicx,booktabs,longtable,array,enumitem}
\usepackage[most]{tcolorbox}
\usepackage{tikz}
\usetikzlibrary{babel,arrows.meta,positioning}
\usepackage{setspace}
\usepackage[hidelinks]{hyperref}
\onehalfspacing
\setlength{\emergencystretch}{3em}
\providecommand{\uadectitulo}{Material de estudio}
\providecommand{\uadecautor}{Material de estudio AulaTeX}
\newcommand{\portadaUAdeC}{%
  \hypersetup{pageanchor=false}
  \begin{titlepage}
  \centering
  {\Large Universidad Autónoma de Coahuila\par}
  \vspace{0.7cm}
  {\large Preparación para ingreso a Médico Cirujano\par}
  \vspace{1.8cm}
  {\LARGE\bfseries\uadectitulo\par}
  \vspace{1cm}
  {\large Actividad 1: síntesis y representación conceptual\par}
  {\large Actividad 2: práctica y soluciones justificadas\par}
  \vfill
  {\uadecautor\par}
  {Revisión: 10 de octubre de 2026\par}
  \vspace{0.8cm}
  \begin{minipage}{0.87\linewidth}
  \small Material formativo local. No constituye una guía, plantilla o examen oficial de UAdeC o CENEVAL. La cobertura y dificultad no están certificadas frente a una convocatoria de ingreso. No sustituye orientación ni atención clínica.
  \end{minipage}
  \end{titlepage}
  \hypersetup{pageanchor=true}
  \tableofcontents\clearpage
  \section*{Alcance y procedencia}
  Las actividades siguen las consignas locales conservadas en esta materia. Los reactivos son de práctica y las respuestas son propuestas razonadas, no resultados de un intento del aspirante. Las fuentes de apoyo se identifican al final. Los PDF bibliográficos locales se inventarían por archivo y hash; su disponibilidad no acredita lectura íntegra ni permiso de redistribución.
}
\begin{document}
\portadaUAdeC
\section{Actividad 1: síntesis de investigación}
\subsection{Adecuación y coherencia}
Adecuación ajusta registro y propósito al destinatario; coherencia organiza ideas sin contradicciones. Un informe debe separar observación de interpretación. No conviene atribuir certeza causal a una asociación.
\subsection{Cohesión y concordancia}
Cohesión relaciona oraciones mediante referentes y conectores. La concordancia nominal relaciona género y número y la verbal, persona y número de sujeto y verbo. Debe distinguirse el núcleo del sujeto de un complemento cercano.
\subsection{Producto de edición}
Borrador original: Los resultados de la encuesta muestra una tendencia. Hubieron veinte respuestas. Sin embargo, por tanto, se requiere revisar la muestra. Versión revisada: Los resultados de la encuesta muestran una tendencia. Hubo veinte respuestas. Por tanto, se requiere revisar la muestra. La primera corrección concuerda con resultados; haber impersonal permanece singular y se elimina el conector adversativo incompatible con la consecuencia.
\subsection{Representación conceptual y cuadro de síntesis}
\begin{longtable}{p{0.27\linewidth}p{0.65\linewidth}}
\toprule Concepto & Relación o función\\\midrule
Propósito & Registro y destinatario\\
Contenido & Coherencia y evidencia\\
Oración & Concordancia y referentes\\
Edición & Corrección sin cambiar sentido\\
\bottomrule
\end{longtable}
\begin{figure}[htbp]
\centering
\begin{tikzpicture}[node distance=5mm, every node/.style={draw,rounded corners=2pt,text width=10cm,align=center,font=\small}]
\node (step0) {Propósito: Registro y destinatario};
\node[below=of step0] (step1) {Contenido: Coherencia y evidencia};
\draw[-{Stealth}] (step0) -- (step1);
\node[below=of step1] (step2) {Oración: Concordancia y referentes};
\draw[-{Stealth}] (step1) -- (step2);
\node[below=of step2] (step3) {Edición: Corrección sin cambiar sentido};
\draw[-{Stealth}] (step2) -- (step3);
\end{tikzpicture}
\caption{Esquema didáctico original; las flechas indican orden de lectura o relación conceptual, no causalidad universal.}
\end{figure}
\clearpage\section{Actividad 2: cuestionario original resuelto}
\begin{enumerate}
\item\textbf{Selecciona la oración con concordancia correcta.}
\begin{enumerate}[label=\Alph*)]
\item Los resultados muestra cambios
\item Las tablas muestran cambios
\item La tabla muestran cambios
\item El resultado muestran cambios
\end{enumerate}
\textbf{Respuesta: B.} Tablas y muestran coinciden en plural; los otros pares no concuerdan.
\item\textbf{Selecciona el uso correcto de haber impersonal.}
\begin{enumerate}[label=\Alph*)]
\item Hubieron tres sesiones
\item Habían una sesión ayer
\item Hubo tres sesiones
\item Han habido tres sesiones
\end{enumerate}
\textbf{Respuesta: C.} Haber existencial es impersonal y se usa en singular: hubo; también ha habido, no han habido.
\item\textbf{¿Qué conector expresa consecuencia?}
\begin{enumerate}[label=\Alph*)]
\item Por tanto
\item Aunque
\item Sin embargo
\item En cambio
\end{enumerate}
\textbf{Respuesta: A.} Por tanto introduce consecuencia; aunque concesión y los otros contraste.
\item\textbf{¿Qué versión distingue evidencia de interpretación?}
\begin{enumerate}[label=\Alph*)]
\item El gráfico demuestra todo
\item Sin datos, es seguro
\item La encuesta prueba causalidad
\item La encuesta sugiere una tendencia en la muestra
\end{enumerate}
\textbf{Respuesta: D.} Sugiere y en la muestra limitan la inferencia; las otras expresiones exageran certeza.
\item\textbf{¿Cuál evita una coma indebida entre sujeto y verbo?}
\begin{enumerate}[label=\Alph*)]
\item La investigación, identifica patrones
\item La investigación identifica patrones
\item La, investigación identifica patrones
\item La investigación identifica, patrones
\end{enumerate}
\textbf{Respuesta: B.} No se separa sujeto de predicado sin inciso; las demás comas rompen la estructura.
\item\textbf{Completa: Las observaciones fueron \_\_\_ por el equipo.}
\begin{enumerate}[label=\Alph*)]
\item revisado
\item revisados
\item revisadas
\item revisada
\end{enumerate}
\textbf{Respuesta: C.} Participio revisadas concuerda con observaciones, femenino plural.
\item\textbf{¿Cuál es una solicitud formal clara?}
\begin{enumerate}[label=\Alph*)]
\item Solicito la revisión del documento adjunto
\item Revísalo, ya sabes cuál
\item Eso que falta, hazlo
\item A ver si lo ven
\end{enumerate}
\textbf{Respuesta: A.} Identifica acción y objeto en registro formal; las otras dependen de referentes vagos.
\item\textbf{Elige la oración con referente inequívoco.}
\begin{enumerate}[label=\Alph*)]
\item Ana habló con Julia cuando ella terminó
\item Lo revisó cuando eso pasó
\item Le dijo que lo hiciera
\item Ana terminó el informe y después habló con Julia
\end{enumerate}
\textbf{Respuesta: D.} La última asigna explícitamente acciones a Ana; las otras tienen referentes ambiguos.
\item\textbf{¿Qué acción mejora coherencia sin inventar contenido?}
\begin{enumerate}[label=\Alph*)]
\item Añadir cifras no medidas
\item Ordenar propósito, método, resultados y límites
\item Eliminar todas las fuentes
\item Sustituir resultados por opiniones
\end{enumerate}
\textbf{Respuesta: B.} Ordenar estructura preserva evidencia; las otras alteran o degradan su respaldo.
\item\textbf{Completa: El conjunto de documentos \_\_\_ disponible.}
\begin{enumerate}[label=\Alph*)]
\item están
\item estuvieron
\item está
\item estaremos
\end{enumerate}
\textbf{Respuesta: C.} El núcleo singular conjunto requiere está; documentos es complemento, no el núcleo.
\end{enumerate}

\section{Conclusiones}
El material permite repasar conceptos y contrastar respuestas. La clave no acredita desempeño del aspirante ni aprobación institucional. Deben revisarse la convocatoria y el temario aplicables antes de usarlo como plan definitivo de ingreso.
\clearpage\section{Fuentes de apoyo y alcance de consulta}
Real Academia Española y ASALE. Diccionario panhispánico de dudas: concordancia, 2.ª edición. Consulta: 10 de octubre de 2026. \url{https://www.rae.es/dpd/concordancia}\par Artículo consultado: concordancia nominal y verbal; ejercicios originales de edición, no reactivos oficiales.\par\medskip
Los originales bibliográficos disponibles se identifican por nombre de archivo y SHA256 en el inventario local. No se atribuyen ediciones, páginas ni lecturas no comprobadas.
\end{document}
